% part: intuitionistic-logic
% chapter: propositions-as-types
% section: reduction

\documentclass[../../../include/open-logic-section]{subfiles}

\begin{document}

\olfileid[hi]{int}{pty}{red}

\olsection{अपचयन}

प्राकृतिक निगमन !!{derivation}s में किसी !!a{introduction} नियम के बाद आने वाला
!!a{elimination} नियम ऐसा चक्कर है जिसका ``अपचयन'' किया जा सकता है। उदाहरणतः
$\Intro{\land}$ के बाद $\Elim{\land}$ आने पर दोनों ``निरस्त'' हो जाते हैं,
क्योंकि $\Elim{\land}$ का निष्कर्ष सदा उस संयोजन के संयोज्यों में से एक है
जिसे \Intro{\land} उत्पन्न करता है, और $\Intro{\land}$ के आधार-वाक्य दोनों
संयोज्य हैं। अतः यदि हमें किसी संयोज्य की !!a{derivation} चाहिए, तो संयोजन को
प्रस्तुत करके फिर उसका विलोपन करने के बजाय सीधे उसी !!{derivation} का उपयोग कर
सकते हैं। हमें मिलता है
\begin{gather*}
  \bottomAlignProof
  \AxiomC{}
  \RightLabel{$\delta_1$}
  \DeduceC{$!A_1$}
  \AxiomC{}
  \RightLabel{$\delta_2$}
  \DeduceC{$!A_2$}
  \RightLabel{$\Intro{\land}$}
  \BinaryInfC{$!A_1 \land !A_2$}
  \RightLabel{$\Elim{\land}_i$}
  \UnaryInfC{$!A_i$}
  \DisplayProof\bottomAlignProof
  \quad
  \redone
  \quad
  \AxiomC{}
  \RightLabel{$\delta_i$}
  \DeduceC{$!A_i$}
  \DisplayProof
\end{gather*}

इस सरलीकरण से सुझाया गया ऐसा ही अपचयन-संबंध संबंधित प्रमाण पदों पर माना जा
सकता है। यदि $N_1$,~$\delta_1$ का और $N_2$,~$\delta_2$ का प्रमाण पद है, तो कह
सकते हैं कि बाएँ प्रमाण का प्रमाण पद (एक चरण में) दाएँ प्रमाण के प्रमाण पद में
अपचयित होता है:
\[
  \proj{i}{\pair{N_1, N_2}} \redone N_i
\]
प्रकारित लैम्ब्डा कलन में यह गुणनफल प्रकार का \emph{बीटा-अपचयन} नियम है।

प्राकृतिक निगमन में \Intro{\land} के बाद \Elim{\land} जैसे, अर्थात् अपचयन के
अधीन, अनुमान-युग्म को \emph{कट} कहते हैं। प्रकारित लैम्ब्डा कलन में
$\redone$ की बाईं ओर के पद को \emph{अपचय्य} और दाईं ओर के पद को
\emph{अपचयफल} कहते हैं। अप्रकारित लैम्ब्डा कलन, जहाँ केवल
$(\lambd[x][N])Q$ को अपचय्य माना जाता है, के विपरीत प्रकारित लैम्ब्डा कलन में
वाक्यविन्यास को $\pair{N}{M}$, $\proj{i}{N}$, $\inj{i}{!A}{N}$,
$\dcase{N}{x_1}{M_1}{x_2}{M_2}$ और $\abort{!A}{N}$ वाले पदों तथा संगत
अपचय्यों तक विस्तृत किया जाता है।

इसी प्रकार विकल्प के लिए अपचयन है:
\begin{gather*}
  \bottomAlignProof
  \AxiomC{}
  \RightLabel{$\delta$}
  \DeduceC{$!A_i$}
  \RightLabel{$\Intro{\lor}$}
  \UnaryInfC{$!A_1 \lor !A_2$}
  \AxiomC{$\Discharge{!A_1}{x_1}$}
  \RightLabel{$\delta_1$}
  \DeduceC{$!C$}
  \AxiomC{$\Discharge{!A_2}{x_2}$}
  \RightLabel{$\delta_2$}
  \DeduceC{$!C$}
  \DischargeRule{$\Elim{\lor}$}{x_1,x_2}
  \TrinaryInfC{$!C$}
  \DisplayProof\bottomAlignProof
  \quad
  \redone
  \quad
  \AxiomC{}
  \RightLabel{$\delta$}
  \DeduceC{$!A_i$}
  \RightLabel{$\delta_i$}
  \DeduceC{$!C$}
  \DisplayProof
\end{gather*}
यह प्रमाण पदों के निम्न अपचयन के अनुरूप है:
\begin{gather*}
  \dcase{\inj[!A_i]{i}{M}}{x_1}{N_1}{x_2}{N_2} \redone \Subst{N_i}{M}{x_i}
\end{gather*}
जहाँ $M$,~$\delta$ का और $N_i$,~$\delta_i$ का प्रमाण पद है। यह योग प्रकारों
का बीटा-अपचयन नियम है। यहाँ $\Subst{M}{N_i}{x}$ का अर्थ $M$ में चर~$x$ की
सभी मुक्त आवृत्तियों को~$N_i$ से बदलना है।

फलन प्रकार का अपचयन $\Intro\lif$ के बाद आने वाले $\Elim\lif$ के चक्कर को
हटाने के अनुरूप है।
\begin{gather*}
  \bottomAlignProof
  \AxiomC{$\Discharge{!A}{x}$}
  \RightLabel{$\delta$}
  \DeduceC{$!B$}
  \DischargeRule{$\Intro{\lif}$}{x}
  \UnaryInfC{$!A \lif !B$}
  \AxiomC{}
  \RightLabel{$\delta'$}
  \DeduceC{$!A$}
  \RightLabel{$\Elim{\lif}$}
  \BinaryInfC{$!B$}
  \DisplayProof\bottomAlignProof
  \quad
  \redone
  \quad
  \AxiomC{}
  \RightLabel{$\delta'$}
  \DeduceC{$!A$}
  \RightLabel{$\delta$}
  \DeduceC{$!B$}
  \DisplayProof
\end{gather*}
प्रमाण पदों के लिए यह साधारण बीटा-अपचयन है:
\begin{gather*}
  (\lambd[\typeof{x}{!A}][N]M)
  \redone \Subst{N}{M}{x}
\end{gather*}

!!{proof}s पर अपचयन रूपांतरणों के अतिरिक्त हम क्रम-विनिमय रूपांतरणों पर विचार
कर सकते हैं। ये ऐसे (उप-)!!{proof}s पर लागू होते हैं जो \Elim{\lor} के बाद
किसी अन्य !!{elimination} नियम पर समाप्त होते हैं, उदाहरणतः:
\begin{multline*}
  \bottomAlignProof
  \AxiomC{}
  \RightLabel{$\delta$}
  \DeduceC{$!A_i$}
  \RightLabel{$\Intro{\lor}$}
  \UnaryInfC{$!A_1 \lor !A_2$}
  \AxiomC{$\Discharge{!A_1}{x_1}$}
  \RightLabel{$\delta_1$}
  \DeduceC{$!C_1 \land !C_2$}
  \AxiomC{$\Discharge{!A_2}{x_2}$}
  \RightLabel{$\delta_2$}
  \DeduceC{$!C_1 \land !C_2$}
  \DischargeRule{$\Elim{\lor}$}{x_1,x_2}
  \TrinaryInfC{$!C_1 \land !C_2$}
  \RightLabel{\Elim{\land}[i]}
  \UnaryInfC{$!C_i$}
  \DisplayProof\bottomAlignProof
  \quad
  \redone
  \\
  \AxiomC{}
  \RightLabel{$\delta$}
  \DeduceC{$!A_i$}
  \RightLabel{$\Intro{\lor}$}
  \UnaryInfC{$!A_1 \lor !A_2$}
  \AxiomC{$\Discharge{!A_1}{x_1}$}
  \RightLabel{$\delta_1$}
  \DeduceC{$!C_1 \land !C_2$}
  \RightLabel{\Elim{\land}[i]}
  \UnaryInfC{$!C_i$}
  \AxiomC{$\Discharge{!A_2}{x_2}$}
  \RightLabel{$\delta_2$}
  \DeduceC{$!C_1 \land !C_2$}
  \RightLabel{\Elim{\land}[i]}
  \UnaryInfC{$!C_i$}
  \DischargeRule{$\Elim{\lor}$}{x_1,x_2}
  \TrinaryInfC{$!C_i$}
  \DisplayProof
\end{multline*}
यदि $\delta$, $\delta_1$ और $\delta_2$ के प्रमाण पद क्रमशः $M$, $N_1$ और
$N_2$ हैं, तो प्रमाण पदों, अर्थात् $\lambd$ पदों, का संगत अपचयन नियम है:
\[
  \proj{i}{\dcase{M}{x_1}{N_1}{x_2}{N_2}} \redone 
\dcase{M}{x_1}{\proj{i}{N_1}}{x_2}{\proj{i}{N_2}}
\]

निश्चय ही हम चक्करों को केवल !!{proof}s के अंत से ही नहीं, बल्कि चक्कर पर
समाप्त होने वाले उप-!!{proof} पर अपचयन रूपांतरण लगाकर !!{proof}s के भीतर से भी
हटा सकते हैं। इसी प्रकार $\beta$-अपचयन, $\lambd$-पदों के उपपदों पर लागू होता
है। यदि $N$, $M$ का उपपद और $N \redone N'$ है, तो~$M$ में उपपद $N$ को $N'$
से बदलकर $M'$ पा सकते हैं। उस स्थिति में $M \redone M'$ भी लिखते हैं। यदि
कोई अनुक्रम $M = M_1 \redone M_2 \redone M_2 \redone \dots \redone M_k =
M'$ है ($k=1$, अर्थात् $M = M'$, वाली स्थिति सहित), तो $M \red M'$ लिखते हैं।

जिस !!{proof} के किसी उप-!!{proof} पर कोई रूपांतरण लागू नहीं किया जा सकता, उसे
\emph{प्रसामान्य} कहते हैं। इसी प्रकार जिस $\lambd$-पद के किसी उपपद पर कोई
रूपांतरण लागू नहीं किया जा सकता, उसे \emph{प्रसामान्य} कहते हैं। प्रसामान्य
$\lambd$-पद को \emph{प्रसामान्य रूप} भी कहते हैं।

\end{document}
